\section{Reflection: The Re-examination of the Evaluation}
\label{p11:sec:reflection}

\noindent
The cycle turns to reflection, the re-examination of the evaluation and of the cognition that framed it. The task of this moment is to hold the structured results of evaluation against the value they proxy and do not reach, and so to prevent the structured measure from being mistaken for the appraisal of the real. The moment draws upon the normative and the virtue-theoretic frameworks, which return here to reassert what the measures cannot contain, and upon the phenomenological framework, which reasserts the lived meaning the structuring discarded. It is the moment at which the operational concession is most stringently enforced, and the moment that prevents the evaluation of \xref{p11:sec:evaluation} from hardening into a possessed truth about the relation.

\subsection{The blindness of the measures, and forged flourishing}
\label{p11:sec:ref-blind}

The first work of reflection is to recall what the measures cannot register. A high score on a satisfaction scale is compatible with a relation that is unjust, since the mean satisfaction is raised by the contentment of a favoured party and is insensitive to the maldistribution beneath it (\xref{p11:sec:ev-skew}); it is compatible with a relation that forecloses one party's unfolding, since the foreclosed party may report satisfaction with a life it has been habituated to expect nothing more from; and it is compatible with a relation whose cycle has settled into a closed repetition that returns nothing new, since the comfort of the familiar registers as satisfaction. A relation may, in short, score well and be a forged flourishing: a satisfaction that has the measurable marks of the good without its reality, in the manner the prior paper analysed as the forged holonomy, the closed cycle that wears the appearance of the spiral.

\begin{claim}[Measured satisfaction is compatible with forged flourishing]
A high measured satisfaction is necessary neither for nor sufficient for a just eudaimonia. It is compatible with injustice, since an aggregate measure is insensitive to the skew of return; with the foreclosure of a party's unfolding, since a foreclosed party may report contentment with a diminished expectation; and with a closed and depleting repetition, since familiarity registers as satisfaction. The measured satisfaction is therefore a structured proxy that a forged flourishing can satisfy, and the work of reflection is to read the measures against the conditions they cannot register, the skew of return, the foreclosure of unfolding, and the closure of the cycle, and so to distinguish a measured satisfaction that marks a real flourishing from one that marks a forged one.
\end{claim}

\noindent
This is the point at which the orthogonal instruments of \xref{p11:sec:evaluation} earn their place. The satisfaction measure alone cannot distinguish the real flourishing from the forged, since both score well on it. The skew of return and the concentration of foreclosed undertakings can distinguish them, since the forged flourishing, however high its mean satisfaction, exhibits the skew and the foreclosure that the real flourishing does not. Reflection is the moment at which the satisfaction measure is read against these orthogonal proxies, and a high satisfaction accompanied by a severe skew is recognised for what it is, the contentment of the favoured party in an unjust relation.

\subsection{The return of the question of justice}
\label{p11:sec:ref-justice}

The second work of reflection is to reassert the question of justice, which the aggregate measures structurally suppress. An evaluation conducted in means and totals asks how satisfying the relation is, and does not ask to whom its satisfactions accrue and at whose cost. Reflection reasserts the second question. It reads the skew of the return distribution as evidence of who is consigned to the low-return tail; it reads the concentration of foreclosed undertakings as evidence of whose unfolding the relation forecloses; and it identifies, where these signatures coincide, the party reduced to the sustaining fuel of the relation, the party who generates value that does not return and whose own unfolding is set aside on account of a shared life that returns to it neither. The question of justice is the question the prior paper held to be constitutive of eudaimonia rather than additional to it, and reflection is the moment at which it is brought to bear on the results of evaluation, against the tendency of the aggregate measures to suppress it.

\subsection{The evaluation may not be possessed as the truth of the relation}
\label{p11:sec:ref-possession}

The third work of reflection is the most consequential, and it is the point at which the operational concession bears most stringently. The structured evaluation of a relation, once conducted, presents a standing temptation: to be taken as the truth of the relation, the authoritative account of how things stand, against which the parties' own appraisals are to be corrected. To yield to this temptation is to commit, in the register of evaluation, the error the prior paper identified as the cardinal one.

\begin{claim}[The evaluation may not be possessed as the truth of the relation]
The structured evaluation of a relation is an operationally warranted proxy, and it may not be possessed as the truth of the relation. To take the measure as the truth, as the authoritative account against which the parties' own appraisals are to be corrected, is to install the evaluation in the place of the barred Other: to treat a structured proxy as the standpoint from which the real value of the relation is known, which is the standpoint the series has established does not exist. The error repeats, in the register of evaluation, the cardinal error of the prior paper, the private certainty of one's own goodness, now as the certainty conferred by a measure; and it is the same occupation of the place of the Other, the same settling of an open and unsymbolised appraisal into a possessed and closed determination. The authority of the evaluation is operational, defeasible, and answerable to the basal appraisal it proxies; it is an instrument of the relation's continued appraisal of itself, and not a verdict upon the relation delivered from outside it. The measure is to be used and not believed, consulted and not obeyed, held as a proxy and never possessed as the truth.
\end{claim}

\noindent
This is the deepest cut of the operational concession. The concession permitted the structured measure, for its operational gain, on the understanding that it is a proxy for a value it does not reach. Reflection enforces the understanding by refusing the measure the authority of truth: the moment the evaluation is possessed as the truth of the relation, the concession is broken, and the proxy that was admitted as an instrument becomes a new metalanguage, a structured idol installed in the place the series has kept empty. The measure that said a relation was well is not the relation's being well, and the measure that said it was ill is not the relation's being ill; both are proxies, to be weighed in the continued appraisal that no measure concludes. Reflection keeps the place of the Other empty against the measure's claim to fill it, and so returns the evaluation to its standing as an instrument of a self-appraisal that remains, in the end, unsymbolised and unconcluded.

\subsection{Reflection as the safeguard of the cycle}
\label{p11:sec:ref-safeguard}

The three works of reflection together safeguard the cycle against the characteristic failures of its evaluative moment. Reading the measures against forged flourishing safeguards against the complacency of a high score; reasserting the question of justice safeguards against the suppression of cost by aggregate; and refusing the measure the authority of truth safeguards against the installation of a structured idol in the place of the real. Reflection is, in this respect, the moment at which the cycle protects itself from its own instruments, and it prepares the moment of re-cognition by returning, corrected and chastened, to the cognition with which the cycle began. The relation that has reflected does not carry forward its measures as truths; it carries forward a revised and re-opened appraisal, in which the measures have been weighed and found to be proxies, and the question of how things really stand between the two has been returned to the unsymbolised register in which it is finally settled, and never finally settled. The cycle now turns to re-cognition, in which this revised appraisal becomes the cognition of a further turn.
